\documentclass[12pt,a4paper,onecolumn]{report}
\usepackage[utf8]{inputenc}
\usepackage{amsmath}
\usepackage{tabularx}
\usepackage{url}
\usepackage{hyperref}


\begin{document}

\noindent
\large \textbf{Kasus Gerak Harmonik Sederhana pada Ayunan} \\
\noindent
\text{Zanuar 23-087, Nabillah 23-092}

\vspace{2em}

\noindent
\section{Latar Belakang} 
Gerak harmonik sederhana (GHS) merupakan dasar penting dalam memahami sistem osilasi \cite{getaranmekanik2020}. 
Selain itu, metode numerik seperti Euler, Heun, dan RK4 digunakan untuk membandingkan akurasi 
dalam menyelesaikan persamaan diferensial yang mendeskripsikan ayunan sederhana. 

\vspace{1.2em}

\section{Solusi Eksak}
Gerak ayunan sederhana dengan sudut kecil dimodelkan oleh persamaan diferensial \cite{getaranmekanik2020}:
$$\frac{d^2x}{dt^2} + \omega^2 x = 0,$$
dengan 

$$\omega = \sqrt{\frac{g}{L}},$$
di mana $g$ adalah percepatan gravitasi dan $L$ panjang tali ayunan.

\vspace{1em}
\noindent
Dengan kondisi awal $x(0) = A$ dan $\dot{x}(0) = 0$, solusi menjadi:
\begin{equation}
x(t) = A \cos(\omega t).
\end{equation}

\noindent
Kecepatan diperoleh \cite{euler_method} dengan turunan pertama dari posisi didapatkan:
\begin{equation}
v(t) = \frac{dx}{dt} = -A \omega \sin(\omega t).
\end{equation}

\newpage
\noindent
dengan:

\begin{itemize}
    \item $x (t)$ = posisi ayunan pada waktu $t$
    \item $v (t)$ = kecepatan ayunan pada waktu $t$
    \item $g = 9.81$ $meter/detik$ = posisi ayunan pada waktu $t$
    \item $L$ = $2$ $meter$ (panjang tali)
    \item $\theta_0 = 0.14$ $rad$ = (Posisi Awal)
    \item $\omega_0 = 0.0$ $rad/detik$ = (Kecepatan Awal)
    \item $h = 0.1$ $detik$ (Langkah)
\end{itemize}

\vspace{2em}


\noindent
\subsection{Euler}
Euler adalah teknik numerik untuk menyelesaikan masalah persamaan diferensial biasa,
metode numerik ini memungkinkan kita untuk meperkirakan solusi untuk persamaan diferensial
menggunakan proses iterative \cite{euler_method}
\vspace{0.3em}

\noindent
\textbf{Metodologi Euler}

\noindent
Metode Euler menggunakan rumus berikut \cite{euler_method_free}:  

\begin{center}
    $y(x + h) = y(x) + h \times f(x, y)$
\end{center}

\subsection{Heun}
Heun adalah

\subsection{Runge Kutta Orde 4}
RK4 adalah

\bibliographystyle{plain}
\bibliography{daftar}

\end{document}
